\documentclass[journal]{IEEEtran}
\usepackage[a5paper, margin=10mm, onecolumn]{geometry}
\usepackage{tfrupee}

\setlength{\headheight}{1cm}
\setlength{\headsep}{0mm}

\usepackage{gvv-book}
\usepackage{gvv}
\usepackage{cite}
\usepackage{amsmath,amssymb,amsfonts,amsthm}
\usepackage{algorithmic}
\usepackage{graphicx}
\usepackage{textcomp}
\usepackage{xcolor}
\usepackage{txfonts}
\usepackage{enumitem}
\usepackage{mathtools}
\usepackage{gensymb}
\usepackage{comment}
\usepackage[breaklinks=true]{hyperref}
\usepackage{tkz-euclide}

\graphicspath{{./figs/}}

\begin{document}
\title{12.296}
\author{AI25BTECH11010 - Dhanush Kumar}
\maketitle
\renewcommand{\thefigure}{\theenumi}
\renewcommand{\thetable}{\theenumi}

\noindent\textbf{Question:}\\
The area enclosed between the straight line $y = x$ and the parabola $y = x^2$ in the $xy$-plane is (PI 2012)

\begin{multicols}{2}
\begin{enumerate}
\item $\dfrac{1}{6}$
\item $\dfrac{1}{4}$
\item $\dfrac{1}{3}$
\item $\dfrac{1}{2}$
\end{enumerate}
\end{multicols}

\noindent\textbf{Solution:}

For $y = x^2$, we write it in matrix form as
\begin{align}
\vec{x}^{\top} V \vec{x} + 2\vec{u}^{\top}\vec{x} + f = 0
\end{align}
where
\begin{align}
V = \myvec{1 & 0 \\ 0 & 0}, \quad 
\vec{u} = \myvec{0 \\ -\tfrac{1}{2}}, \quad 
f = 0.
\end{align}

The line $y = x$ can be written parametrically as
\begin{align}
\vec{x} = \vec{h} + \kappa \vec{m}, \quad \kappa \in \mathbb{R}
\end{align}
where
\begin{align}
\vec{h} = \myvec{0 \\ 0}, \quad 
\vec{m} = \myvec{1 \\ 1}.
\end{align}

The intersection points satisfy
\begin{align}
\kappa_i =
\frac{-\vec{m}^{\top}(V\vec{h} + \vec{u}) 
\pm \sqrt{[\vec{m}^{\top}(V\vec{h} + \vec{u})]^2 
- (\vec{m}^{\top}V\vec{m})\,g(\vec{h})}}
{\vec{m}^{\top}V\vec{m}}
\end{align}
where
\begin{align}
g(\vec{h}) = \vec{h}^{\top}V\vec{h} + 2\vec{u}^{\top}\vec{h} + f.
\end{align}

Substituting values:
\begin{align}
V\vec{h} + \vec{u} &= \myvec{0 \\ -\tfrac{1}{2}}, \\
\vec{m}^{\top}(V\vec{h} + \vec{u}) &= -\tfrac{1}{2}, \\
\vec{m}^{\top}V\vec{m} &= 1, \\
g(\vec{h}) &= 0.
\end{align}

Hence,
\begin{align}
\kappa_{1,2} &= 
\frac{-(-\tfrac{1}{2}) \pm \sqrt{(-\tfrac{1}{2})^2 - 1(0)}}{1} \\
&= \frac{\tfrac{1}{2} \pm \tfrac{1}{2}}{1}.
\end{align}

Therefore,
\begin{align}
\kappa_1 = 0, \quad \kappa_2 = 1.
\end{align}

The intersection points are
\begin{align}
\vec{x}_1 = \vec{h} + \kappa_1 \vec{m} = \myvec{0 \\ 0}, \quad
\vec{x}_2 = \vec{h} + \kappa_2 \vec{m} = \myvec{1 \\ 1}.
\end{align}

The area enclosed between the line and the parabola is
\begin{align}
A &= \int_{0}^{1} (y_{\text{line}} - y_{\text{parabola}})\,dx \\
  &= \int_{0}^{1} (x - x^2)\,dx \\
  &= \left[\frac{x^2}{2} - \frac{x^3}{3}\right]_0^1 \\
  &= \frac{1}{2} - \frac{1}{3} \\
  &= \frac{1}{6}.
\end{align}

\begin{align}
\boxed{A = \frac{1}{6}}
\end{align}

\noindent\textbf{Final Answer: (a) } $\dfrac{1}{6}$

\end{document}

